% TOP_PROBLEM: {"problemId": "problem.square-peg-problem-toeplitz-s-inscribed-square-conjecture", "canonicalTitle": "Square Peg Problem (Toeplitz's Inscribed Square Conjecture)", "releaseRank": 489, "primaryDomain": "geometry_topology", "primaryDomainLabel": "Geometry & topology", "displayStatus": "open_with_solved_subcases", "releaseStatus": "open_with_solved_subcases"}
% !TeX program = lualatex
% ATTEMPT_STATUS: unresolved
\documentclass[11pt]{article}
\usepackage[margin=25mm]{geometry}
\usepackage{fontspec}
\setmainfont{Latin Modern Roman}
\usepackage{amsmath,amssymb,mathtools}
\usepackage{unicode-math}
\setmathfont{Latin Modern Math}
\usepackage{xurl,hyperref}
\hypersetup{colorlinks=true,urlcolor=blue}
\setcounter{secnumdepth}{0}
\setlength{\parindent}{0pt}
\setlength{\parskip}{5pt}
\setlength{\emergencystretch}{3em}
\begin{document}
\section*{\texorpdfstring{Square Peg Problem (Toeplitz's Inscribed Square Conjecture)}{Square Peg Problem (Toeplitz's Inscribed Square Conjecture)}}\label{489}
\subsection{Definitions and mathematical statement}
A Jordan curve is the image $\Gamma=\gamma(S^1)\subset{\mathbb{R}}^2$ of a continuous injective map from the circle. Let $J$ be rotation through $\pi/2$. The square-peg assertion is
\[
 \exists x,v\in{\mathbb{R}}^2,\ v\ne0:\quad
 x,\ x+v,\ x+v+Jv,\ x+Jv\in\Gamma.
\]
These are four distinct vertices of a Euclidean square. The curve need not be differentiable, rectifiable, or convex, and the square need not lie inside its bounded complementary region. Only the vertices are required to lie on the curve. A limiting configuration whose side length tends to zero does not supply a nondegenerate square.
\par\smallskip\textit{Formulation and concept references:} \href{https://arxiv.org/pdf/2203.02613}{[S1]}.

\subsection{Short English statement}
Does every continuous simple closed plane curve contain the four vertices of a nondegenerate square?
\subsection{Research attempt}
\paragraph{Scope and literature check.}
This attempt by GPT-6 Astra Ultra was prepared on 27 September 2026. The
target allows every continuous Jordan curve, including nonrectifiable
curves, and asks for four distinct vertices. No assertion that a square's
interior lies inside the curve is added. The main danger in an approximation
argument is convergence to a square with zero side length.

The source [S1], Theorem 1.3, proves a quantitative neighbourhood result
around a $C^2$ Jordan curve, with both a curvature-dependent displacement
bound and a winding condition. Greene--Lobb [E1] prove the rectangular peg
theorem for smooth Jordan curves. The January 2026 version 3 of Asano--Ike
[E2], Theorem 1.1 and Corollaries 1.2--1.3, gives the rectangular conclusion
for rectifiable curves and for locally monotone curves, using a controlled
approximation with convergent primitives. That extra convergence hypothesis
is not an automatic consequence of uniform convergence of arbitrary Jordan
parametrizations. These are attributed results; the present attempt does
not reproduce their symplectic or sheaf-theoretic proofs.

The plan is to prove an elementary symmetric subcase, isolate a separate
positive-area subcase, and then attack the limiting obstruction directly.
The last route produces precise sufficient conditions and a counterexample
to an overly strong compactness shortcut, but not the missing condition
for an arbitrary zero-area Jordan curve.

\paragraph{Lemma 489.1: centrally symmetric radial curves have a square.}
Suppose, after translation, that
\[
 \Gamma=\{r(\theta)(\cos\theta,\sin\theta):\theta\in\mathbb R\},
 \qquad r>0,\quad r\text{ continuous},\quad
 r(\theta+\pi)=r(\theta).
\]
This describes a Jordan curve: the radius is positive and two distinct
directions modulo $2\pi$ cannot give the same point. Let
\[
                      f(\theta)=r(\theta)-r(\theta+\pi/2).
\]
The symmetry gives $f(\theta+\pi/2)=-f(\theta)$. If $f(0)=0$ we
already have a zero; otherwise the intermediate value theorem on
$[0,\pi/2]$ produces one. At a zero $\theta_*$, all four radii in
the directions $\theta_*+j\pi/2$ are equal to some $r_*>0$.
The resulting four points are the vertices of a square centred at the
origin, with side length $\sqrt2 r_*$. In particular
\begin{equation}\label{a489:radiallower}
                    \text{side length}\ge\sqrt2\min_\theta r(\theta)>0.
\end{equation}
Every centrally symmetric convex body with interior has such a radial
description about its centre, so this also proves that familiar convex
subcase. Neither convexity nor differentiability was needed for the stated
radial lemma. The argument is a rederivation of an elementary special case,
not a new solution of the general conjecture.

\paragraph{Test of the symmetry-removal step.}
A tempting extension is to look for four equal quarter-turn radii for an
arbitrary star-shaped curve. It fails even for
\begin{equation}\label{a489:asymmetric}
                         r(\theta)=2+\tfrac1{10}\cos\theta.
\end{equation}
The radius is everywhere positive, so this is a smooth embedded radial
curve. If its opposite radii at $\theta$ and $\theta+\pi$ were equal,
we would have $\cos\theta=0$. Equality of the other opposite pair would
also require $\sin\theta=0$, which is impossible. Thus this curve has
no square \emph{centred at the chosen origin}. It is not a counterexample
to the square-peg conjecture: the centre of the required square is free,
and the smooth theorem applies to this curve.

This is a useful diagnostic of the proof rather than merely a missing
line. Central symmetry supplied two of the three radius equalities and
turned the last one into a sign-changing scalar function. Without it,
there is no reason a single angular parameter solves all three equations.
Allowing the centre to vary is natural, but the curve need not admit a
single-valued radial parametrization from every prospective centre.

\paragraph{Lemma 489.2: positive planar measure supplies squares directly.}
Let $E\subset\mathbb R^2$ be measurable with $0<|E|<\infty$.
Then for every sufficiently small nonzero vector $v$, there exists $x$ with
\begin{equation}\label{a489:measure-square}
                         x,x+v,x+Jv,x+v+Jv\in E.
\end{equation}
\emph{Proof.} Translation is continuous in $L^1(\mathbb R^2)$:
\[
          \|\mathbf{1}_E(\cdot+w)-\mathbf{1}_E\|_1\longrightarrow0
                                  \quad\text{as }w\to0.
\]
One elementary justification approximates $\mathbf{1}_E$ in $L^1$ by
finite linear combinations of indicators of rectangles. Translation
continuity for a rectangle follows by bounding the area of the strips along
its sides. Translation preserves the $L^1$ norm, so the approximation
error contributes at most twice its original size. Sending that error and
then $|w|$ to zero proves the assertion.

Use the convention $E-w=\{x:x+w\in E\}$ and set
\[
             I_v=E\cap(E-v)\cap(E-Jv)\cap(E-v-Jv).
\]
The union bound, with no probabilistic independence assumption, gives
\[
 |I_v|\ge |E|-|E\setminus(E-v)|-|E\setminus(E-Jv)|
                           -|E\setminus(E-v-Jv)|.
\]
All three subtracted quantities tend to zero. For small enough $|v|$,
their sum is less than $|E|$, so $I_v$ is nonempty. Any $x\in I_v$
satisfies \eqref{a489:measure-square}. Since $v\ne0$, the four points
are distinct and form a square.\hfill$\square$

Apply this to $E=\Gamma$ whenever the Jordan curve itself has positive
two-dimensional measure. Boundedness gives finite measure, and the lemma
settles that subcase even without rectifiability or local monotonicity.
This positive-measure observation is also recorded in [E2], Remark 5.5;
the proof above makes its measure-theoretic mechanism explicit. The area
of the \emph{enclosed region} cannot be used in its place. Every Jordan
curve encloses a region of positive area, but applying the lemma to that
region places the square's vertices in the region, not on its boundary.
That substitution would solve a different and much easier problem.

\paragraph{Attempt A: encode all squares without an order assumption.}
Write $z_i=\gamma(t_i)$ and consider the continuous map from $(S^1)^4$
to $\mathbb R^4$ given by
\begin{equation}\label{a489:testmap}
 \Phi(t_1,t_2,t_3,t_4)=
 \left(z_1+z_3-z_2-z_4,\;
       |z_1-z_3|^2-|z_2-z_4|^2,\;
       (z_1-z_3)\cdot(z_2-z_4)\right).
\end{equation}
The first component is a vector in $\mathbb R^2$, and the other two
are scalars. A zero with $z_1\ne z_3$ is exactly a nondegenerate square,
with $z_1,z_3$ its opposite pair. Indeed the first equation gives a common
midpoint $c$. Write the pairs as $c\pm a$ and $c\pm b$; the other
equations say $|a|=|b|$ and $a\cdot b=0$. Since $a\ne0$, also $b\ne0$,
and these are precisely the four vertices of a square. Conversely a square
clearly gives such a zero. If $z_1=z_3$ at a zero, equal diagonal lengths
and the midpoint equation force all four points to coincide.

Thus the target is an off-diagonal zero of \eqref{a489:testmap}. The map
has a whole circle of diagonal zeros $t_1=t_2=t_3=t_4$, regardless of the
curve. Counting variables and equations, applying compactness, or finding
the minimum value zero of $|\Phi|$ cannot distinguish those forced zeros
from the ones wanted. No unproved degree computation is inserted here.
For numerical searches one must explicitly exclude a neighbourhood of the
diagonal or enforce a positive separation bound, and then prove that the
excluded neighbourhood does not carry all relevant zeros.

\paragraph{Lemma 489.3: a quantitative condition that makes passage to the limit valid.}
Suppose compact sets $\Gamma_j$ converge in Hausdorff distance to a
compact set $\Gamma$, and $\Gamma_j$ contains square vertices
$x_j,x_j+v_j,x_j+v_j+Jv_j,x_j+Jv_j$ with $|v_j|\ge c>0$.
Since Hausdorff convergence makes the sets uniformly bounded, a subsequence
has $x_j\to x$ and $v_j\to v$. The limiting points lie in $\Gamma$:
their distance to $\Gamma$ tends to zero and $\Gamma$ is closed.
The square relations pass to the limit and $|v|\ge c$, proving that
$\Gamma$ contains a nondegenerate square.

There is an equivalent useful sufficient condition in parameter space.
Let $\gamma_j\to\gamma$ uniformly, with $\gamma$ a Jordan parametrization.
For $0<\delta\le\pi$ put
\[
 m_\delta=\min\{|\gamma(s)-\gamma(t)|:
                   d_{S^1}(s,t)\ge\delta\}>0.
\]
Positivity follows from injectivity and compactness of the set of separated
parameter pairs. If some pair of the four chosen square parameters has
circle distance at least $\delta$, and
$\epsilon_j=\|\gamma_j-\gamma\|_\infty$, then its vertex distance is
at least $m_\delta-2\epsilon_j$. The diameter of a square is $\sqrt2$
times its side, hence
\begin{equation}\label{a489:noncollapse}
                 |v_j|\ge\frac{m_\delta-2\epsilon_j}{\sqrt2}.
\end{equation}
For all large $j$ this is a fixed positive lower bound. What is missing
in the unrestricted approximation attack is the selection of squares with
such a separation, uniformly along an approximating sequence.

\paragraph{Attempt B: obtain separation from local geometry.}
No graph $y=f(x)$ with Lipschitz constant $L<1$ contains four square
vertices. To see this, take the two perpendicular sides issuing from one
vertex and write their displacement vectors as $(a,b)$ and $(c,d)$.
Graph injectivity gives $a,c\ne0$, and the Lipschitz bound gives
$|b|\le L|a|$ and $|d|\le L|c|$. Perpendicularity would imply
\[
                     |ac|=|bd|\le L^2|ac|<|ac|,
\]
a contradiction. The assertion is unchanged under rotation of the graph.

Consequently, if a family of curves has a common $\eta>0$ such that
every subset of diameter less than $\eta$ is contained in a local
rotated graph of Lipschitz constant less than 1, every inscribed square
in that family has side at least $\eta/\sqrt2$. The limiting lemma now
applies. For an individual regular $C^1$ embedded curve, compactness and
continuity of its nonzero tangent provide finitely many graph charts with
slopes bounded by some $L<1$ and a positive Lebesgue number for this cover.
This explains the exclusion of collapsed squares at that level of
regularity. It does not give a uniform Lebesgue number for arbitrary smooth
curves merely converging in $C^0$.

Here is an explicit test of the latter failure. Start with the boundary
of the unit square. For $0<\delta<1/4$, replace its bottom boundary segment
from $(1/2,0)$ to $(1/2+\delta,0)$ by the three segments joining
\[
 (1/2,0),\quad(1/2,\delta),\quad
 (1/2+\delta,\delta),\quad(1/2+\delta,0).
\]
The resulting polygonal curve is simple: the three inserted segments lie
inside the original square and meet the unchanged boundary only at their
two endpoints. It converges in Hausdorff distance to the unit-square
boundary, with distance at most $\delta$. Its four inserted corner points
form an inscribed square of side $\delta$. The chosen squares therefore
collapse although the limit is still a Jordan curve.

The same approximants also have a large inscribed square, namely their
four unchanged outer corners. Thus this construction disproves the claim
that \emph{every choice} of squares in converging Jordan curves is
noncollapsing. It does not disprove the existence of a better selection,
which is exactly what a valid approximation proof would need. Nor is it
a counterexample to the conjecture itself.

\paragraph{Verification and next steps.}
The companion script checks the midpoint, perpendicularity, and equal-length
conditions for rationally parametrized rotated squares, and verifies the
notched polygons and their tiny squares for several rational $\delta$.
It also numerically locates sign changes in centrally symmetric radial
examples. These computations stress-test the formulas and the counterexample
to an intermediate claim; the existence statements are justified by the
proofs above. A grid search can miss a square on a continuous curve, and
an approximate numerical zero of \eqref{a489:testmap} does not certify an
exact off-diagonal zero.

\textbf{UNRESOLVED.} The proved elementary subcases are centrally symmetric
positive radial curves and all curves having positive planar measure.
The noncollapse lemma gives a rigorous sufficient condition for extending
an approximation argument, while the notch construction shows why an
arbitrary selection of squares is inadequate. The first remaining gap is
to produce a uniformly separated square along suitable approximants of
an arbitrary Jordan curve outside the established classes, or to bypass
that selection altogether. Three precise directions are to find an invariant
forcing an off-diagonal zero of \eqref{a489:testmap}, to control the primitives
required by [E2] for a larger class of zero-area curves, or to prove a lower
bound for the largest square in an approximant from data that survive
uniform convergence. None of those general statements has been proved here.


\subsection{Sources}
\begin{itemize}
\item[S1]Gregory R. Chambers, "On the Square Peg Problem," arXiv:2203.02613 [math.GT], March 2022.
\url{https://arxiv.org/pdf/2203.02613}.
\textit{Formulation source.}
\item[E1]Joshua Evan Greene and Andrew Lobb,
\emph{The rectangular peg problem}, Annals of Mathematics 194 (2021),
509--517; author manuscript arXiv:2005.09193v1.
\url{https://arxiv.org/abs/2005.09193v1}.
\textit{Attributed smooth-curve theorem; consulted 27 September 2026.}
\item[E2]Tomohiro Asano and Yuichi Ike,
\emph{The rectifiable rectangular peg problem}, arXiv:2412.21057v3,
5 January 2026, Theorem 1.1, Corollaries 1.2--1.3, and Remark 5.5.
\url{https://arxiv.org/html/2412.21057v3}.
\textit{Attributed manuscript results with their approximation hypotheses;
the positive-measure subcase is independently proved in this attempt.
Consulted 27 September 2026.}
\end{itemize}
\end{document}
